\section{Fundamentação Tecnológica: Nanopartículas e Dispersoides em RHEAs}\label{sec:fundamentacao_nanoparticulas_rhea}

O desenvolvimento de Ligas de Alta Entropia Refratárias (RHEAs) reforçadas com nanopartículas e dispersoides representa a fronteira atual na engenharia de materiais para ambientes aerotérmicos severos, tais como bordos de ataque, cones de nariz, escudos térmicos e sistemas de propulsão \emph{scramjet}. A incorporação sinérgica de fases nanométricas visa superar o clássico conflito de projeto metalúrgico (\emph{trade-off}) entre resistência mecânica em ultra-alta temperatura (acima de $1500~^\circ\mathrm{C}$) e ductilidade à temperatura ambiente.

\subsection{Sistemas de Matrizes e Nanoreforços na Literatura}\label{subsec:sistemas_matrizes_nanoreforcos}

A literatura recente mapeada organiza-se primordialmente em torno de matrizes refratárias de rede cúbica de corpo centrado (BCC) baseadas em elementos dos Grupos IV, V e VI (tais como NbMoTaW, NbMoTaWHf, WTaV, TiVNbTa e HfNbTaTiZr), combinadas a dispersoides cerâmicos avançados. Dentre os principais reforços explorados, destacam-se os nitretos de alta entropia (HENs) e de háfnio ($\mathrm{HfN}$ FCC e $\mathrm{MN}$ HCP) dispersos em matrizes BCC, a exemplo do sistema $\mathrm{NbMoTaW(HfN)}_x$, bem como carbetos de ultra-alta temperatura (UHTCs) mono e multimetálicos ($\mathrm{TaC}$, $\mathrm{HfC}$, $\mathrm{ZrC}$, $\mathrm{TiC}$ e o carbeto de alta entropia $\mathrm{TiZrHfNbTaC}$) integrados via rotas \emph{in-situ} ou \emph{ex-situ}.

Adicionalmente, investiga-se a formação de heteroestruturas BCC-cerâmica com boretos refratários do tipo $M_3\mathrm{B}_2$ ($M = \mathrm{Ta}, \mathrm{V}, \mathrm{W}$) dotados de alta coesão interfacial. Na classe das ligas com óxidos dispersos (ODS-HEAs), a adição de óxidos termodinamicamente ultraestáveis ($\mathrm{Y}_2\mathrm{O}_3$, $\mathrm{ZrO}_2$ e $\mathrm{HfO}_2$) viabilizada por técnicas de não-equilíbrio --- como a Ablação por Varredura a Laser (LSA) --- possibilita a síntese de nanopartículas híbridas com diâmetros rigorosamente controlados entre $11~\mathrm{nm}$ e $56~\mathrm{nm}$. No domínio dos silicetos, nanoestruturas do tipo $M_5\mathrm{Si}_3$ ($400$--$600~\mathrm{nm}$) e revestimentos baseados em $(\mathrm{NbMoTaW})\mathrm{Si}_2$ dopados com ítrio operam como barreiras térmicas autorregenerativas.

\subsection{Mecanismos de Atuação sob Condições Extremas ($> 1500~^\circ\mathrm{C}$)}\label{subsec:mecanismos_atuacao_extremas}

A estabilização mecânica e a resistência à deformação por fluência sob calor extremo são governadas primordialmente por dois fenômenos físicos: o Mecanismo de Orowan, no qual partículas nanométricas dispersas impõem o contorno forçado (\emph{bypass}) das discordâncias móveis ($\tau = Gb/\lambda$), elevando de modo expressivo o limite de escoamento em regimes térmicos elevados; e a Ancoragem de Zener (\emph{Zener Pinning}), na qual dispersoides localizados em contornos de grão e junções triplas restringem a difusão atômica e suprimem o crescimento celular excessivo. Na liga $\mathrm{NbMoTaWHfN}$, por exemplo, a presença da fase $\mathrm{HfN}$ retém o tamanho de grão em aproximadamente $5{,}5~\mu\mathrm{m}$ (em contraste com $\sim 118~\mu\mathrm{m}$ na liga não reforçada), preservando o endurecimento Hall-Petch até $1800~^\circ\mathrm{C}$. Paralelamente, em sistemas como $(\mathrm{WTaV})_{90}\mathrm{B}_{10}$, a segregação controlada de boro origina contornos de fase com ordenamento gradiente (GOPBs), mitigando concentrações localizadas de tensão e propiciando cerca de 6\% de plasticidade à temperatura ambiente.

No que tange à degradação termoquímica e ablação aerotérmica, os nanoreforços aceleram a passivação através do fornecimento de sítios preferenciais para a formação rápida e contínua de óxidos termodinamicamente densos ($\alpha\text{-}\mathrm{Al}_2\mathrm{O}_3$, $\mathrm{Cr}_2\mathrm{O}_3$, $\mathrm{SiO}_2$ e vidros borossilicatos autorregenerativos). Em contrapartida, RHEAs monolíticas desprovidas desse controle microestrutural sofrem desintegração catastrófica por \emph{pest oxidation} entre $600~^\circ\mathrm{C}$ e $1000~^\circ\mathrm{C}$ devido à expansão volumétrica desestabilizadora de óxidos voláteis de molibdênio e tungstênio.

\subsection{Desafios Termodinâmicos e Estabilidade de Fases}\label{subsec:desafios_termodinamicos_limitacoes}

Apesar das propriedades promissoras, a operação em temperaturas superiores a $0{,}6\,T_m$ ($> 1500~^\circ\mathrm{C}$) impõe desafios cinéticos expressivos. O amadurecimento de Ostwald (\emph{Ostwald ripening}) induz a dissolução gradual de nanopartículas menores em favor de agregados maiores, ampliando o espaçamento interpartícula ($\lambda$) e degradando a eficiência do ancoramento mecânico. Podem ocorrer ainda reações interfaciais deletérias com precipitação de compostos frágeis, como fases de Laves intermetálicas ($\mathrm{Cr}_2\mathrm{Nb}$) e segregação de impurezas em contornos de grão. Por fim, frações volumétricas excessivas de dispersoides cerâmicos elevam substancialmente a Temperatura de Transição Dúctil-Frágil (DBTT), tornando indispensável o controle estequiométrico e a dopagem precisa de intersticiais (como nitrogênio e boro) para preservar a integridade mecânica a frio.

\subsection{Panorama Quantitativo de Desempenho (Síntese da Literatura)}\label{subsec:panorama_quantitativo_desempenho}

A Tabela~\ref{tab:desempenho_nanoparticulas_rhea} consolida as principais formulações de RHEAs nano-reforçadas e seus respectivos regimes operacionais sob alta temperatura reportados na literatura.

\begin{table}[!htbp]
\centering
\caption{Panorama quantitativo de desempenho mecânico e térmico de sistemas RHEA reforçados com dispersoides e nanopartículas.}\label{tab:desempenho_nanoparticulas_rhea}
\footnotesize
\begin{tabularx}{\textwidth}{>{\raggedright\arraybackslash}p{2.2cm} >{\raggedright\arraybackslash}p{2.3cm} >{\raggedright\arraybackslash}p{2.0cm} >{\raggedright\arraybackslash}p{1.8cm} >{\raggedright\arraybackslash}X}
\toprule
\textbf{Sistema / Liga} & \textbf{Fase de Reforço / Dispersoide} & \textbf{Limite de Escoamento ($\sigma_{0{,}2}$)} & \textbf{Temperatura de Teste} & \textbf{Principais Desempenhos e Métricas} \\
\midrule
\textbf{NbMoTaWHfN} & $\mathrm{HfN}$ (FCC) + $\mathrm{MN}$ (HCP) & $1682~\mathrm{MPa}$\newline $792~\mathrm{MPa}$\newline $288~\mathrm{MPa}$ & $25~^\circ\mathrm{C}$\newline $1400~^\circ\mathrm{C}$\newline $1800~^\circ\mathrm{C}$ & Deformação $> 30\%$ a $1800~^\circ\mathrm{C}$; amolecimento estimado em $1543~^\circ\mathrm{C}$. \\
\midrule
\textbf{$(\mathrm{WTaV})_{90}\mathrm{B}_{10}$} & Boreto $M_3\mathrm{B}_2$ + GOPBs (B) & $650~\mathrm{MPa}$\newline $242~\mathrm{MPa}$ & $25~^\circ\mathrm{C}$\newline $1600~^\circ\mathrm{C}$ & Plasticidade $\sim 6\%$ à temperatura ambiente; estabilidade até $\sim 0{,}7\,T_m$. \\
\midrule
\textbf{$\mathrm{TiVNbTaSi}_{0{,}1}$} & Nano-silicetos ($M_5\mathrm{Si}_3$) & Alta resistência ao amolecimento & Até $1000~^\circ\mathrm{C}$ & Mantém ductilidade em tração à temperatura ambiente sem fases frágeis após $50~\mathrm{h}$. \\
\midrule
\textbf{$(\mathrm{NbMoTaW})\mathrm{Si}_2$ + Y} & Silicieto modificado & Barreira de alta integridade & Até $2100~^\circ\mathrm{C}$ & Camada vítrea de sílica autorregenerativa (\emph{self-healing}) contra ablação. \\
\midrule
\textbf{$\mathrm{Hf}_{0{,}5}\mathrm{Nb}_{0{,}5}\mathrm{Ta}_{0{,}5}\mathrm{Ti}_{1{,}5}\mathrm{Zr}$} & Monolítica (sem nano-reforço) & Falha por \emph{Pest Oxidation} & $600$--$1000~^\circ\mathrm{C}$ & Variação mássica severa ($-6{,}022~\mathrm{mg/mm}^2$ a $800~^\circ\mathrm{C}$) por esfarelamento. \\
\bottomrule
\end{tabularx}
\legend{Fonte: Elaborado pelos autores (2026), com base em dados compilados da literatura.}
\end{table}
